%% kapa: exam=1 round=24 year=2013 period=1 form=A subject=civil_law subject_ko=민법 seq=33 number=33
%% source: Q-Net 감정평가사 기출문제 (cst00309) artlSeq=5073802 file=extracted/1008/5073802/2013년 제24회 감정평가사 1차시험 기출문제/1교시(A형).hwp
%% conversion: kordoc 4.24.51 -> tools/gen1.py (auto); review=proofread-sonnet; flags=none
\begin{question}{33}
지역권에 관한 설명으로 옳은 것은? (다툼이 있으면 판례에 의함)
\begin{choices}
  \item 하나의 승역지에 순차로 여러 개의 용수지역권이 설정된 경우 지역권 상호간에는 우열이 없다.
  \item 요역지 공유자의 1인에 대해 지역권의 소멸시효 정지사유가 있어도 다른 공유자의 지역권의 소멸시효가 정지되지는 않는다.
  \item 일정한 장소를 오랜 시일 통행한 사실이 있다면 통로의 개설이 없더라도 지역권을 시효취득할 수 있다.
  \item 요역지를 제3자가 시효취득하면 제3자는 그 토지를 위한 지역권도 취득한다.
  \item 요역지는 1필의 토지일 필요는 없으나 승역지는 1필의 토지이어야 한다.
\end{choices}
\end{question}
\answer{4}
